% ------------------------------------------------------------
% Chapter 2: Comparison with Other Emergent Gravity Theories
% ------------------------------------------------------------

\chapter{Comparison with Other Emergent Gravity Theories}
\label{chap:comparison-emergent}

\section{Overview}

A central claim of the lamphrodynamic framework is that gravitational
phenomena emerge from entropy gradients and thermodynamic fluxes in the
plenum. This places \rsvpabbr{} in conceptual proximity to several other
emergent--gravity proposals, most prominently those of Jacobson,
Verlinde, and Padmanabhan, as well as more recent quantum--information
approaches (notably Carney). The goal of this chapter is to compare these
theories in detail, highlighting overlaps, divergences, and potential
experimental discriminants.

Throughout, we work at the level of qualitative mechanism, then translate
to quantitative predictions where possible. A comparison table summarizing
assumptions and key predictions appears in \cref{sec:comparison-table}.

\section{Jacobson's Thermodynamic Gravity}
\label{sec:jacobson}

\subsection{Local Rindler Horizons and Clausius Relation}

Jacobson's 1995 derivation of Einstein's equations interprets gravity as a
thermodynamic equation of state. Consider a local Rindler horizon $H$
generated by accelerated observers. The Clausius relation
\[
  \delta Q = T\,\mathrm{d}S
\]
applied to local horizons yields, under certain assumptions,
\[
  G_{\mu\nu} + \Lambda g_{\mu\nu} = 8\pi G\,T_{\mu\nu},
\]
where $G_{\mu\nu}$ is the Einstein tensor and $T_{\mu\nu}$ the stress--energy
tensor of matter fields. Crucially, this derivation presupposes a specific
relation between area and entropy ($\mathrm{d}S \propto \mathrm{d}A$), borrowed
from black--hole thermodynamics.

\subsection{Comparison with Lamphrodynamic Entropy}

\rsvpabbr{} shares the theme that entropy plays a gravitational role, but
differs in three fundamental ways:
\begin{enumerate}
\item Entropy in \rsvpabbr{} is \emph{volume--based} rather than
horizon--based: $\SField$ is defined locally throughout the plenum, without
reference to causal horizons.
\item Gravity arises from \emph{gradients and fluxes} of $\SField$ rather
than equilibrium relations at horizons.
\item The field equations are modified Einstein equations with an explicit
lamphrodynamic contribution $\Theta_{\mu\nu}$ derived from the AKSZ action, not
Einstein's equations recovered from thermodynamic identities.
\end{enumerate}

Thus, whereas Jacobson infers Einstein's equations from thermodynamic
assumptions, \rsvpabbr{} derives \emph{modified} Einstein equations from a
microscopic lamphrodynamic field theory.

\section{Verlinde's Entropic Gravity}
\label{sec:verlinde}

\subsection{Holographic Screens and Entropic Force}

Verlinde proposes that gravity is an entropic force arising from changes of
entropy associated with holographic screens. A test mass approaching a screen
experiences an entropic force
\[
  F = T\,\frac{\partial S}{\partial x},
\]
leading, in heuristic derivations, to Newtonian gravity. Later developments
attempt to explain dark--matter--like phenomena through entropic effects
associated with de Sitter entropy.

\subsection{Differences with \rsvpabbr{}}

Although both frameworks employ entropy gradients, several sharp differences
exist:
\begin{enumerate}
\item Verlinde's entropy is associated with \emph{holographic screens} and
fundamentally tied to information storage capacity; \rsvpabbr{} entropy is a
local thermodynamic field defined in the bulk.
\item Verlinde's derivation is essentially Newtonian or weak--field, whereas
\rsvpabbr{} yields fully covariant modified Einstein equations.
\item Verlinde's approach postulates emergent space from holography;
\rsvpabbr{} assumes a lamphrodynamic plenum with intrinsic geometry, not
constructed from entanglement entropy.
\end{enumerate}

Consequently, \rsvpabbr{} predicts distinct strong--field and cosmological
deviations from general relativity, while Verlinde's modifications primarily
affect galaxy--scale dynamics.

\section{Padmanabhan's Emergent Paradigm}
\label{sec:padmanabhan}

\subsection{Degrees of Freedom Balance}
Padmanabhan argues that spacetime dynamics emerge from the thermodynamic
equipartition of degrees of freedom between bulk and boundary, encoded in
holographic relations such as
\[
  N_{\mathrm{sur}} - N_{\mathrm{bulk}} \sim \text{cosmological acceleration}.
\]
Acceleration of cosmic expansion results from a mismatch between horizon and
bulk degrees of freedom.

\subsection{Comparison with \rsvpabbr{}}

The lamphrodynamic theory shares with Padmanabhan the idea that cosmological
acceleration is thermodynamic in origin. However:
\begin{enumerate}
\item \rsvpabbr{} attributes acceleration to entropy \emph{flux} and gradient
terms in $\Theta_{\mu\nu}$, not mismatch of horizon degrees of freedom.
\item Cosmology in \rsvpabbr{} does not rely on de Sitter entropy or horizon
thermodynamics.
\item \rsvpabbr{} does not posit holographic equipartition as a fundamental
principle; instead, AKSZ variations yield modified field equations directly.
\end{enumerate}

\section{Quantum Information Approaches (Carney)}
\label{sec:carney}

\subsection{Gravity from Decoherence and Entanglement}

Carney and collaborators propose that gravity may emerge from quantum
information dynamics such as decoherence, entanglement structure, or
information flow in quantum systems. These approaches focus on microscopic
quantum mechanisms rather than macroscopic thermodynamics.

\subsection{Relevance to \rsvpabbr{}}

Lamphrodynamic entropy $\SField$ is classical and thermodynamic rather than
quantum informational. However, a deep connection may exist if $\SField$
coarse--grains microscopic quantum correlations. If so, lamphrodynamic fluxes
could represent emergent semiclassical remnants of underlying quantum
information flows. This is speculative but suggests potential connections
between AKSZ quantization and holographic quantum information geometry.

\section{Comparison Table}
\label{sec:comparison-table}

\begin{center}
\begin{tabular}{l|c|c|c|c}
\textbf{Property} & \textbf{RSVP} & \textbf{Jacobson} &
\textbf{Verlinde} & \textbf{Padmanabhan} \\
\hline
Entropy type &
volume, local &
horizon &
holographic screen &
horizon/bulk balance \\
\hline
Mechanism &
gradients, flux &
Clausius relation &
entropic force &
DOF mismatch \\
\hline
Field equations &
modified Einstein &
Einstein &
Newtonian/heuristic &
cosmological dynamical law \\
\hline
Dark--matter effects &
entropy flux &
not primary &
entropic de Sitter &
holographic acceleration \\
\hline
Quantum origin &
AKSZ, classical limit &
thermo heuristic &
holography heuristic &
holographic DOF
\end{tabular}
\end{center}

\section{Summary and Outlook}

While \rsvpabbr{} is thermodynamic in spirit, it differs from other
emergent--gravity proposals in mechanism, ontology, and mathematical
structure. The lamphrodynamic corrections arise from explicit field
equations, not from thermodynamic identities or holographic assumptions. In
the following chapter, we turn to observable consequences for geodesic
motion and test particles, including perihelion precession and light bending,
which provide immediate empirical tests of the lamphrodynamic corrections.

\clearpage

